% Clean, formal report-style document class and packages
\documentclass[11pt, twoside, a4paper]{article}
\usepackage[utf8]{inputenc}
\usepackage[a4paper, margin=2.5cm]{geometry}

% Essential packages for a formal report
\usepackage{mathtools, amsmath, amssymb, derivative}    % Math environments
\usepackage{float, graphicx, booktabs, subcaption, array} % Figures and tables
\usepackage{listings}                                    % Source code
\usepackage{xcolor}                                     % Colors
\usepackage{titlesec}                                   % Section formatting
\usepackage{pdfpages}                                   % PDF inclusion
\usepackage[style=apa]{biblatex}                        % Bibliography
\usepackage[hidelinks]{hyperref}                        % Links
\usepackage{siunitx}                                    % SI units
\usepackage[version=4]{mhchem}                          % Chemical formulae
\usepackage{enumitem}                                   % List environments
\usepackage{menukeys}                                   % Menu formatting
\usepackage{tcolorbox}                                  % Boxed environments

% Bibliography resources
\addbibresource{references.bib}
\addbibresource{sample.bib}

% Color definitions for code listings
\definecolor{codegreen}{rgb}{0,0.6,0}
\definecolor{codegray}{rgb}{0.5,0.5,0.5}
\definecolor{codepurple}{rgb}{0.58,0,0.82}
\definecolor{backcolour}{rgb}{1,1,1}

% Code listing style
\lstdefinestyle{mystyle}{
    backgroundcolor=\color{backcolour},   
    commentstyle=\color{codegreen},
    keywordstyle=\color{magenta},
    numberstyle=\tiny\color{codegray},
    stringstyle=\color{codepurple},
    basicstyle=\ttfamily\footnotesize,
    breakatwhitespace=false,         
    breaklines=true,                 
    captionpos=b,                    
    keepspaces=true,                 
    numbers=left,                    
    numbersep=5pt,                  
    showspaces=false,                
    showstringspaces=false,
    showtabs=false,                  
    tabsize=2,
    frame=topbottom
}
\lstset{style=mystyle}

% Custom box for outputs
\newtcolorbox{outputbox}[1][]{
    colback=white,
    title=Output,
    #1
}

% Math operators
\DeclareMathOperator{\proj}{proj}

% Custom column types
\newcommand*{\carry}[1][1]{\overset{#1}}
\newcolumntype{B}[1]{r*{#1}{@{\,}r}}

% Formal report formatting
\setlength{\parindent}{0pt}                             % No paragraph indentation
\setlength{\parskip}{0.8em}                             % Space between paragraphs

% Section spacing for formal reports
\titlespacing*{\section}{0pt}{2.5ex plus 1ex minus 0.2ex}{1.3ex plus 0.2ex}
\titlespacing*{\subsection}{0pt}{2.25ex plus 1ex minus 0.2ex}{1ex plus 0.2ex}
\titlespacing*{\subsubsection}{0pt}{1.5ex plus 1ex minus 0.2ex}{0.75ex plus 0.1ex}

% List formatting for formal reports
\setlist[itemize]{leftmargin=*, itemsep=0.5em, topsep=0.5em}
\setlist[enumerate]{leftmargin=*, itemsep=0.5em, topsep=0.5em}
\setlist[itemize]{label=\textbullet}                    % Use bullets for lists
\setlist[itemize,2]{label=\textendash}                  % En-dash for second level

% Hide underfull warnings
\raggedbottom

% Custom theorem environments
\newtheorem{definition}{Definition}

% Custom commands for document information
\newcommand{\TitleName}{Monitoring Toilet Usage with Sensors}
\newcommand{\unitTime}{Semester 1, 2026}
\newcommand{\unitName}{EGH419: Advanced Design and Entrepreneurship}
\newcommand{\documentAuthors}{Xanthea Vazey: n11525762}

\begin{document}
% -------------------------------------------------
%  Title Page
% -------------------------------------------------
\begin{titlepage}
    \thispagestyle{empty}
    \vspace*{0.2\textheight}
    
    \noindent
    \begin{minipage}[t]{0.03\textwidth}
        \rule[-25ex]{1pt}{140pt}   % longer vertical line
    \end{minipage}
    \hspace{1em}
    \begin{minipage}[t]{0.9\textwidth}
        \raggedright
        \normalsize{\unitTime} \\[2ex]
        \Huge\textbf{\TitleName} \\[2ex]
        \normalsize{\unitName}
    \end{minipage}
    
    \vfill
    \begin{flushleft}
        \documentAuthors
    \end{flushleft}
\end{titlepage}

% Page numbering setup
\pagenumbering{roman}
\newpage
\pagenumbering{arabic}

% -------------------------------------------------
%  Main Content
% -------------------------------------------------


\section{Problem Statement}
\subsection{Problem Statement}

Remote sanitation infrastructure is an ongoing operational challenge for land managers responsible for national parks and remote recreation areas. Increasing visitation to protected areas places additional pressure on existing toilet facilities, particularly those located in remote environments where conventional sewer infrastructure is unavailable.

\subsubsection{Problem Drivers}

\begin{itemize}
    \item National parks receive millions of visitors annually, resulting in significant demand for sanitation infrastructure.
    \item Many toilet facilities in remote recreation areas are ``drop toilets'' or pit systems that operate without connection to sewer or wastewater treatment networks.
    \item Waste removal and servicing of these toilets often requires specialised logistics such as long-distance transport or helicopter access, significantly increasing operational costs.
    \item Park managers typically rely on scheduled inspections or manual reporting to determine when toilets require servicing, making it difficult to accurately monitor when systems are approaching capacity.
\end{itemize}

\subsubsection{Supporting Evidence}

\begin{itemize}
    \item Remote regions of the Australian Alps National Parks contain more than 200 toilet installations servicing hiking and camping areas.
    \item Waste from high-use remote toilets may require helicopter extraction when pits or storage systems reach capacity.
    \item Poorly managed human waste in outdoor recreation environments can contribute to soil contamination, water pollution, and potential disease transmission.
\end{itemize}

\subsubsection{Engineering Opportunity}

These challenges present an opportunity for the development of a low-power remote monitoring system capable of measuring toilet fill levels and transmitting status data to park managers. Such a system would allow maintenance teams to service facilities only when required, reducing operational costs, preventing overflow events, and improving environmental protection in remote recreation areas.
\section{Resources}
\begin{itemize}
    \item \url{https://pmc.ncbi.nlm.nih.gov/articles/PMC7586839/}
    \item \url{https://www.nepc.gov.au/sites/default/files/2022-09/regional-remote-case-studies.pdf}
    \item \url{https://www.environment.nsw.gov.au/topics/parks-reserves-and-protected-areas/park-management/npws-park-management/pollution-monitoring}
    \item \url{https://parks.tas.gov.au/Documents/Appendix%20D%20-%20Onsite%20Wastewater%20Assessment%20%28Lake%20Adelaide%29.pdf}
    \item \url{https://www.abc.net.au/news/2023-12-24/waterless-outback-toilet-trial-sanitation-tourism-possibilities/103225172}
    \item \url{https://theaustralianalpsnationalparks.org/wp-content/uploads/2023/08/ALPS-eblast-94.pdf}
    \item \url{https://theaustralianalpsnationalparks.org/wp-content/uploads/2017/07/toilet-report-2001.pdf}
    \item toilet map: \url{https://www.continence.org.au/national-public-toilet-map?utm_source=chatgpt.com} around 23000 toilets aus wide \url{https://toiletmap.gov.au/about}
    \subitem but i think that is all toilets over aus... maybe 
    \item over 200 toilets in the australian alps national parks \url{https://theaustralianalpsnationalparks.org/wp-content/uploads/2017/07/toilet-report-2001.pdf} which will need pumping out or replacement when full 
    \item in NZ there are over 2000 toilets in remote areas \url{https://www.doc.govt.nz/globalassets/documents/getting-involved/volunteer/walking-track-maintenance/walking-track-maintenance-toilet-report.pdf}

\end{itemize}

\section{Similar Products}
\begin{itemize}
    \item Ultrasonic smart bin monitoring system
    \item example company : Enevo or bigbelly
\end{itemize}
Technology used:

ultrasonic fill-level sensors, IoT telemetry (LoRa, NB-IoT, satellite) ,predictive waste collection
These systems measure how full a bin is and optimise servicing.

\subsection{Other things to note}
current methods of getting waste out is helicopter, good to document this and how they know when to change it. then look at cost estimates etc.. 
\begin{itemize}
    \item The strongest justification for your project would be numbers like:
    \item how often remote toilets are serviced
    \item cost of helicopter waste removal
    \item examples of toilets overflowing or being closed
    \item documented helicopter waste removal operations
    \item maintenance cost estimates
    \item visitor-to-toilet ratios in national parks
\end{itemize}


\end{document}